\subsection{মহাকাশ অভিযান}

\begin{learningbox}
এই অংশ শেষে শিক্ষার্থী পারবে:
\begin{itemize}
\item মহাকাশ অভিযানে ICT-এর ভূমিকা ব্যাখ্যা করতে।
\item স্যাটেলাইট, GPS, remote sensing ও weather forecasting-এর ব্যবহার লিখতে।
\item মহাকাশ গবেষণায় data communication ও computer control-এর প্রয়োজনীয়তা বুঝতে।
\end{itemize}
\end{learningbox}

\subsubsection{মহাকাশ অভিযানে আইসিটির ভূমিকা}

মহাকাশ অভিযান হলো satellite, spacecraft, rover, telescope বা astronaut mission-এর মাধ্যমে মহাকাশ সম্পর্কে তথ্য সংগ্রহ ও গবেষণা। এখানে ICT ছাড়া mission control, communication, navigation, data processing, image analysis, simulation এবং safety monitoring সম্ভব নয়।

\begin{definitionbox}{মহাকাশ অভিযানে ICT}
মহাকাশ অভিযানে তথ্য ও যোগাযোগ প্রযুক্তি হলো computer, sensor, communication network, satellite link, software ও data analysis system-এর ব্যবহার, যার মাধ্যমে spacecraft নিয়ন্ত্রণ, তথ্য আদান-প্রদান ও গবেষণা করা হয়।
\end{definitionbox}

\subsubsection{স্যাটেলাইট যোগাযোগ}

Satellite পৃথিবীর কক্ষপথে থেকে communication signal গ্রহণ ও প্রেরণ করে। TV broadcast, mobile backhaul, internet connectivity, disaster communication এবং remote area communication-এ satellite গুরুত্বপূর্ণ।

\begin{examplebox}{বাংলাদেশ প্রেক্ষিত}
Bangabandhu Satellite-1 বাংলাদেশের broadcasting, telecommunication ও emergency communication সক্ষমতা বাড়ানোর জন্য গুরুত্বপূর্ণ national space infrastructure।
\end{examplebox}

\subsubsection{রিমোট সেন্সিং}

Remote sensing হলো দূর থেকে sensor ব্যবহার করে পৃথিবীর land, water, forest, crop, weather বা disaster সম্পর্কে তথ্য সংগ্রহ করা। Satellite image বিশ্লেষণ করে বন্যা, নদীভাঙন, বন উজাড়, ফসলের অবস্থা ও নগরায়ণ পর্যবেক্ষণ করা যায়।

\begin{keypointbox}{রিমোট সেন্সিংয়ের ব্যবহার}
\begin{itemize}
\item দুর্যোগ ব্যবস্থাপনা: বন্যা, ঘূর্ণিঝড়, আগুন।
\item কৃষি: crop health, irrigation planning।
\item পরিবেশ: forest monitoring, pollution observation।
\item নগর পরিকল্পনা: land use map, road network।
\item জলবায়ু গবেষণা: temperature, cloud, sea level trend।
\end{itemize}
\end{keypointbox}

\subsubsection{জিপিএস}

GPS বা Global Positioning System satellite signal ব্যবহার করে অবস্থান, গতি ও সময় নির্ণয় করে। Navigation, transport tracking, ride sharing, disaster rescue, mapping, agriculture এবং military operation-এ GPS ব্যবহার হয়।

\subsubsection{আবহাওয়া পূর্বাভাস}

Weather satellite cloud pattern, temperature, humidity, wind movement ও storm formation পর্যবেক্ষণ করে। Computer model এই data বিশ্লেষণ করে forecast তৈরি করে। ঘূর্ণিঝড়প্রবণ বাংলাদেশে satellite-based weather forecasting জীবন ও সম্পদ রক্ষায় সহায়ক।

\subsubsection{মহাকাশ গবেষণায় তথ্য বিশ্লেষণ}

Space mission থেকে বিপুল পরিমাণ image, sensor data, telemetry ও scientific measurement আসে। Computer algorithm, AI, high-performance computing এবং data visualization ব্যবহার করে বিজ্ঞানীরা এই data বিশ্লেষণ করেন।

\begin{center}
\begin{tabular}{|p{0.30\textwidth}|p{0.56\textwidth}|}
\hline
\textbf{ICT component} & \textbf{মহাকাশে ভূমিকা} \\
\hline
Sensor & temperature, radiation, image, motion data সংগ্রহ \\
\hline
Communication link & spacecraft ও mission control-এর মধ্যে data আদান-প্রদান \\
\hline
Software & navigation, control, simulation, fault detection \\
\hline
Data analysis & satellite image, telemetry ও research data বিশ্লেষণ \\
\hline
\end{tabular}
\end{center}

\begin{boardbox}{বোর্ড ফোকাস}
মহাকাশ অভিযানে ICT-এর ভূমিকা লিখতে satellite communication, remote sensing, GPS, weather forecasting এবং data analysis—এই পাঁচটি point খুব গুরুত্বপূর্ণ।
\end{boardbox}

\begin{practicebox}
\textbf{MCQ}
\begin{enumerate}
\item দূর থেকে satellite sensor দিয়ে পৃথিবীর তথ্য সংগ্রহকে কী বলে?\\
ক. Remote sensing \quad খ. Printing \quad গ. Formatting \quad ঘ. Encryption only
\item GPS প্রধানত কী নির্ণয় করে?\\
ক. file size \quad খ. অবস্থান ও সময় \quad গ. font style \quad ঘ. browser history
\end{enumerate}

\textbf{সংক্ষিপ্ত প্রশ্ন}
\begin{enumerate}
\item মহাকাশ অভিযানে satellite communication-এর ভূমিকা লেখ।
\item আবহাওয়া পূর্বাভাসে ICT কীভাবে সহায়তা করে?
\end{enumerate}
\end{practicebox}

\begin{sourcebox}
এই অংশে local ICT PDF resources, NASA educational material, NOAA/NASA satellite weather information, GPS.gov overview এবং Bangladesh satellite context মিলিয়ে লেখা হয়েছে।
\end{sourcebox}
